\paragrafos{Συμμετρία ως προς σημείο: μισή στροφή}{Β.2.4, «Συμμετρία ως προς σημείο»}

\begin{anagke}
Στο πάρκο δίπλα στο σχολείο υπάρχει μια ρόδα λούνα παρκ. Η Αγγελική κοιτάζει δύο καμπίνες που βρίσκονται η μία απέναντι από την άλλη, και παρατηρεί: «Όταν η ρόδα κάνει μισή στροφή, η πρώτη καμπίνα πηγαίνει ακριβώς στη θέση της δεύτερης!». Ο Νίκος ρωτάει: «Και η δεύτερη στη θέση της πρώτης». Το κέντρο της ρόδας είναι ακριβώς στη μέση ανάμεσα στις δύο καμπίνες.

Η Πολυξένη τα αντιγράφει: «Αν ένα σημείο $\mathrm{Α}$ περιστραφεί γύρω από το κέντρο $\mathrm{Ο}$ κατά $180^\circ$, πηγαίνει στο σημείο $\mathrm{Α}'$ με το $\mathrm{Ο}$ στη μέση του $\mathrm{ΑΑ}'$». Η Μυρτώ το δοκιμάζει με ένα διαφανές χαρτί και μια καρφίτσα: καρφώνει στο $\mathrm{Ο}$, σχεδιάζει ένα τρίγωνο, γυρίζει το χαρτί μισή στροφή και σχεδιάζει το νέο. «Βγήκε ανάποδο, αλλά ίσο!».

Σε αυτή την παράγραφο ορίζουμε τη συμμετρία ως προς σημείο και βλέπουμε πώς ταιριάζει με τη συμμετρία ως προς άξονα. Θα δούμε και ότι οι αντίθετοι αριθμοί στον άξονα είναι συμμετρικά σημεία ως προς το $0$.
\end{anagke}

\begin{lexi}{συμμετρικό ως προς σημείο}{κέντρο συμμετρίας, συμμετρικό}
Το \lexiword{κέντρο συμμετρίας} είναι το σημείο γύρω από το οποίο περιστρέφουμε. Το \lexiword{συμμετρικό} του $\mathrm{Α}$ ως προς το $\mathrm{Ο}$ είναι το σημείο στην «απέναντι» θέση, σε ίση απόσταση από το $\mathrm{Ο}$ και στην ίδια ευθεία.
\end{lexi}

\begin{metafrasi}
Μισή στροφή γύρω από το $\mathrm{Ο}$. & στροφή $180^\circ$\\
Το $\mathrm{Ο}$ είναι το μέσο του $\mathrm{ΑΑ}'$. & $\mathrm{ΟΑ}=\mathrm{ΟΑ}'$ και $\mathrm{Α},\mathrm{Ο},\mathrm{Α}'$ συνευθειακά\\
Το συμμετρικό του $\mathrm{Ο}$ ως προς το $\mathrm{Ο}$. & το $\mathrm{Ο}$\\
Οι αντίθετοι αριθμοί. & συμμετρικά ως προς το $0$\\
\end{metafrasi}

\tmima{Η θεωρία}

\begin{orismos}{συμμετρικό σημείου ως προς σημείο}
\textbf{Συμμετρικό} του σημείου $\mathrm{Α}$ ως προς σημείο $\mathrm{Ο}$ είναι το σημείο $\mathrm{Α}'$ για το οποίο το $\mathrm{Ο}$ είναι το \textbf{μέσο} του τμήματος $\mathrm{ΑΑ}'$. Το $\mathrm{Ο}$ λέγεται \textbf{κέντρο συμμετρίας}. Το συμμετρικό του $\mathrm{Ο}$ είναι το ίδιο το $\mathrm{Ο}$.
\end{orismos}

Με άλλα λόγια, το $\mathrm{Α}'$ προκύπτει αν περιστρέψουμε το $\mathrm{Α}$ κατά $180^\circ$ γύρω από το $\mathrm{Ο}$. Για να το κατασκευάσω: ενώνω το $\mathrm{Α}$ με το $\mathrm{Ο}$, προεκτείνω την $\mathrm{ΑΟ}$ και με τον διαβήτη παίρνω $\mathrm{ΟΑ}'=\mathrm{ΟΑ}$.

\begin{center}
\begin{tikzpicture}[scale=1.1]
\fill (0,0) circle (2pt) node[above right] {$\mathrm{Ο}$};
\fill (-1.8,1.1) circle (2pt) node[above left] {$\mathrm{Α}$}; \fill (1.8,-1.1) circle (2pt) node[below right] {$\mathrm{Α}'$};
\draw[thick] (-1.8,1.1)--(1.8,-1.1);
\fill (-2.2,-0.4) circle (2pt) node[left] {$\mathrm{Β}$}; \fill (2.2,0.4) circle (2pt) node[right] {$\mathrm{Β}'$};
\draw[thick] (-2.2,-0.4)--(2.2,0.4);
\draw[gblue,thick] (-1.8,1.1)--(-2.2,-0.4); \draw[gblue,thick] (1.8,-1.1)--(2.2,0.4);
\end{tikzpicture}
\end{center}

\begin{kanonas}{συμμετρία ως προς σημείο}
Δύο σχήματα είναι \textbf{συμμετρικά ως προς σημείο} $\mathrm{Ο}$ αν το καθένα αποτελείται από τα συμμετρικά των σημείων του άλλου. Τα συμμετρικά ως προς σημείο σχήματα είναι \textbf{ίσα}. Το συμμετρικό τμήματος $\mathrm{ΑΒ}$ είναι τμήμα $\mathrm{Α}'\mathrm{Β}'$ ίσο και \textbf{παράλληλο} προς το $\mathrm{ΑΒ}$.
\end{kanonas}

Στον άξονα των ρητών, το συμμετρικό του αριθμού $\alpha$ ως προς το $0$ είναι ο $-\alpha$, ο αντίθετός του. Ως προς το $1$: το συμμετρικό του $3{,}5$ είναι το $x$ με $\frac{3{,}5+x}{2}=1$, δηλαδή $x=-1{,}5$.

\begin{prosoxi}
\begin{itemize}[leftmargin=1.2em,itemsep=0.3ex,topsep=2pt]
\item Στη συμμετρία ως προς σημείο, τα $\mathrm{Α}$, $\mathrm{Ο}$, $\mathrm{Α}'$ είναι πάνω στην ίδια ευθεία (στη συμμετρία ως προς άξονα το $\mathrm{ΑΑ}'$ είναι κάθετο στον άξονα).
\item Το $\mathrm{Ο}$ είναι το μέσο του $\mathrm{ΑΑ}'$: $\mathrm{ΑΑ}'=2\mathrm{ΟΑ}$.
\item Δεν είναι ίδιο με τη συμμετρία ως προς άξονα: μια στροφή $180^\circ$ δεν «αναποδογυρίζει» το σχήμα σαν καθρέφτης.
\item Το συμμετρικό τμήματος είναι \emph{παράλληλο} στο αρχικό (ή στην ίδια ευθεία, αν η ευθεία περνά από το $\mathrm{Ο}$).
\end{itemize}
\end{prosoxi}

\tmima{Λυμένα παραδείγματα}

\paradeigma{Η ρόδα}
\tixero Η καμπίνα $\mathrm{Α}$ απέχει $12$ μ. από το κέντρο $\mathrm{Ο}$ της ρόδας.\par
\tithelo Πού είναι η συμμετρική καμπίνα και πόσο απέχουν οι δύο.\par
\possindeo Το $\mathrm{Α}'$ είναι στην προέκταση της $\mathrm{ΑΟ}$, με $\mathrm{ΟΑ}'=12$ μ. Το $\mathrm{ΑΑ}'=2\cdot12=24$ μ. (διάμετρος της ρόδας).\par
\apantisi $\mathrm{ΟΑ}'=12$ μ. και $\mathrm{ΑΑ}'=24$ μ.\par
\elegxos Το $\mathrm{Ο}$ είναι το μέσο του $\mathrm{ΑΑ}'$.

\paradeigma{Στον άξονα των αριθμών}
\tixero Το σημείο $\mathrm{Α}(3{,}5)$ του άξονα.\par
\tithelo Το συμμετρικό του ως προς το $0$ και ως προς το $1$.\par
\possindeo Ως προς το $0$: ο αντίθετος $-3{,}5$. Ως προς το $1$: το $1$ είναι το μέσο, $\frac{3{,}5+x}{2}=1$, $x=-1{,}5$.\par
\apantisi $-3{,}5$ και $-1{,}5$.\par
\elegxos Η απόσταση του $3{,}5$ από το $1$ είναι $2{,}5$ και του $-1{,}5$ από το $1$ επίσης $2{,}5$.

\paradeigma{Συμμετρικό τμήματος}
\tixero Τμήμα $\mathrm{ΑΒ}=4$ cm και το συμμετρικό του $\mathrm{Α}'\mathrm{Β}'$ ως προς $\mathrm{Ο}$.\par
\tithelo Το μήκος του και η θέση του.\par
\possindeo Τα συμμετρικά σχήματα είναι ίσα, άρα $\mathrm{Α}'\mathrm{Β}'=4$ cm, και το $\mathrm{Α}'\mathrm{Β}'$ είναι παράλληλο προς το $\mathrm{ΑΒ}$.\par
\apantisi $\mathrm{Α}'\mathrm{Β}'=4$ cm και $\mathrm{Α}'\mathrm{Β}'\parallel\mathrm{ΑΒ}$.\par
\elegxos Το $\mathrm{ΑΒΒ}'\mathrm{Α}'$ έχει τα $\mathrm{Α}\mathrm{Α}'$ και $\mathrm{ΒΒ}'$ να περνούν από το $\mathrm{Ο}$.

\paradeigma{Κατασκευή}
\tixero Το τρίγωνο $\mathrm{ΑΒΓ}$ και το σημείο $\mathrm{Ο}$ εκτός αυτού.\par
\tithelo Να κατασκευάσω το συμμετρικό του.\par
\possindeo Βρίσκω τα $\mathrm{Α}'$, $\mathrm{Β}'$, $\Gamma'$ (το καθένα με $\mathrm{ΟΑ}'=\mathrm{ΟΑ}$ κ.λπ. στην προέκταση) και ενώνω.\par
\apantisi Το τρίγωνο $\mathrm{Α}'\mathrm{Β}'\Gamma'$ είναι ίσο με το $\mathrm{ΑΒΓ}$.\par
\elegxos Με διαφανές χαρτί και στροφή $180^\circ$ γύρω από το $\mathrm{Ο}$ το $\mathrm{ΑΒΓ}$ πέφτει πάνω στο $\mathrm{Α}'\mathrm{Β}'\Gamma'$.

\tmima{Ασκήσεις}

\zoni{Το κατάλαβα}{σημεία, τμήματα, σχήματα}

\ex Σχεδίασε σημείο $\mathrm{Ο}$ και σημείο $\mathrm{Α}$ και βρες το συμμετρικό του ως προς $\mathrm{Ο}$ με διαβήτη.

\ex Αν $\mathrm{ΟΑ}=3{,}2$ cm, πόσο είναι το $\mathrm{ΟΑ}'$ και πόσο το $\mathrm{ΑΑ}'$;

\ex Σχεδίασε τμήμα $\mathrm{ΑΒ}$ και το συμμετρικό του ως προς σημείο $\mathrm{Ο}$ εκτός του. Τι σχέση έχουν;

\ex Σχεδίασε τρίγωνο $\mathrm{ΑΒΓ}$ και το συμμετρικό του ως προς ένα σημείο $\mathrm{Ο}$.

\ex Βρες το συμμετρικό των σημείων του άξονα $\mathrm{Α}(3{,}5)$, $\mathrm{Β}(-2)$, $\Gamma(0)$ ως προς την αρχή $0$.

\ex Βρες το συμμετρικό του $3{,}5$ ως προς το $1$ και του $-2$ ως προς το $3$.

\ex Σωστό ή λάθος; α) Στη συμμετρία ως προς σημείο, τα $\mathrm{Α}$, $\mathrm{Ο}$, $\mathrm{Α}'$ είναι συνευθειακά. \quad β) Το $\mathrm{Ο}$ είναι το μέσο του $\mathrm{ΑΑ}'$. \quad γ) Τα συμμετρικά σχήματα ως προς σημείο είναι ίσα. \quad δ) Το συμμετρικό τμήματος είναι κάθετο σε αυτό.

\zoni{Το χρησιμοποιώ}{ρόδες, ρολόγια, παρτέρια}

\ex \textbf{Η ρόδα.} Μια καμπίνα απέχει $12$ μ. από το κέντρο της ρόδας. Πού είναι η καμπίνα μετά από μισή στροφή; Πόσο απέχουν οι δύο θέσεις;

\ex \textbf{Τα δύο παρτέρια.} Δύο ίδια παρτέρια του κήπου είναι συμμετρικά ως προς το συντριβάνι $\mathrm{Ο}$. Ένα φυτό στο πρώτο απέχει $2{,}5$ μ. από το $\mathrm{Ο}$. Πού είναι το αντίστοιχο στο δεύτερο;

\ex \textbf{Το ρολόι.} Ο λεπτοδείκτης στις $12{:}00$ δείχνει προς τα πάνω. Πού δείχνει μετά από $30$ λεπτά (μισή στροφή); Τι σχέση έχουν οι δύο θέσεις της άκρης του;

\ex \textbf{Ο ανεμόμυλος.} Ένας ανεμόμυλος με δύο αντίθετα πτερύγια περιστρέφεται. Πού βρίσκεται το ένα πτερύγιο όταν το άλλο έχει μετακινηθεί στη θέση του αντιθέτου του;

\ex \textbf{Το χαρτί.} Κάρφωσε με καρφίτσα ένα διαφανές χαρτί, σχεδίασε μια λέξη και γύρνα το $180^\circ$. Πώς φαίνεται η λέξη;

\zoni{Το σκέφτομαι}{ίδιο ή διαφορετικό από τον καθρέφτη;}

\ex Σε τι διαφέρει η συμμετρία ως προς σημείο από τη συμμετρία ως προς άξονα στο συμμετρικό ενός τριγώνου;

\ex Αν το $\mathrm{Α}'$ είναι συμμετρικό του $\mathrm{Α}$ ως προς $\mathrm{Ο}$, ποιο είναι το συμμετρικό του $\mathrm{Α}'$ ως προς το ίδιο $\mathrm{Ο}$;

\ex Εξήγησε γιατί το συμμετρικό ενός τμήματος ως προς σημείο είναι παράλληλο προς αυτό. (Βοήθεια: ένωσε τα $\mathrm{Α}$, $\mathrm{Β}'$ και δες το $\mathrm{ΑΒΑ}'\mathrm{Β}'$.)

\ex Ποια είναι η ευθεία που είναι συμμετρική του εαυτού της ως προς σημείο $\mathrm{Ο}$;

\zoni{Για το σπίτι}{στροφή μισή}

\ex Σχεδίασε ένα σχήμα σε διαφανές χαρτί, κάρφωσε ένα σημείο και γύρνα μισή στροφή. Σχεδίασε το συμμετρικό.

\ex Βρες στο σπίτι δύο αντικείμενα που ταιριάζουν αν περιστραφούν $180^\circ$.
